% Source-aligned English for printed p.120.
% Decisions: TOH-P120-R1-EN-001 through TOH-P120-R1-EN-008;
% TOH-P120-R2-REV-001 adds the zero-width footnote semantic target.
\SourcePage{120}

\LayerTarget{tohoku:introduction:I:p4:cont:p120}
the theory also applies to the non-Hausdorff spaces that occur in
abstract algebraic geometry or in ``arithmetic geometry'' [15] [8].
Conversations with R. Godement and H. Cartan were very valuable in working out
the theory, and in particular Godement's introduction of \emph{flasque sheaves}
and \emph{soft sheaves}, which advantageously replace fine sheaves in many
questions, proved extremely convenient. A more complete account, to which we
shall refer for various points of detail, will be given in a book in
preparation by R. Godement [9].

\LayerTarget{tohoku:introduction:I:p5}
Chapter IV treats the nonclassical question of Ext for sheaves of modules; in
particular, one finds there a useful spectral sequence relating ``global'' Ext
and ``local'' Ext.
\LayerTarget{tohoku:introduction:I:p6}
The situation becomes more complicated in Chapter V, where, in addition, a
group $G$ acts on the space $X$, the sheaf of rings $\mathbf{O}$ given on $X$, and the
sheaves of $\mathbf{O}$-modules under consideration. In particular, in 5.2 one obtains
a statement that seems to me to be the definitive form of the ``Čechist''
cohomological theory of spaces with a (non-topological) group of operators,
which may have fixed points. It is expressed by introducing new functors
$H^n(X;G,A)$ (already implicit in many earlier special cases): one then finds
two spectral functors, with remarkable initial terms, which abut to them.

\LayerTarget{tohoku:introduction:II}
\LayerContents{II. Applications}
\paragraph{II. Applications.}
\LayerTarget{tohoku:introduction:II:p1}
For lack of space, I have been able to give in this paper only very few
applications of the techniques used (notably in 3.4 and 3.6), contenting
myself with pointing out a few of them in passing. Let us also mention the
following applications:

\LayerTarget{tohoku:introduction:II:a}
\par a)
The notion of Ext for sheaves of modules permits the most general known
formulation of Serre's ``algebraic duality theorem'': If $A$ is a coherent
algebraic sheaf [15] on a nonsingular projective algebraic variety of dimension
$n$, then the dual of $H^p(X,A)$ is canonically identified with
$\mathrm{Ext}_{\mathrm{O}}^{n-p}(X;A,\Omega^n)$, where $\mathbf{O}$ (respectively $\Omega^n$) is
the sheaf of germs of regular functions (respectively regular $n$-forms) on
$X$.

\LayerTarget{tohoku:introduction:II:b}
\par b)
All the formalism developed in Chapters III, IV, V can be applied in abstract
algebraic geometry. I shall show elsewhere how it permits one to extend to
complete algebraic varieties various results proved by Serre [15] [16] [17]
for projective varieties.

\LayerTarget{tohoku:introduction:II:c}
\par c)
It seems that the $H^n(X;G,A)$ are the natural intermediary for a general
theory of reduced Steenrod powers in sheaves, and of the cohomology of
symmetric powers of arbitrary spaces, a theory which also applies in
algebraic geometry in characteristic $p$.

\LayerTarget{tohoku:introduction:III}
\LayerContents{III. Gaps}
\paragraph{III. Gaps.}
\LayerTarget{tohoku:introduction:III:p1}
To avoid lengthening this paper, I have passed over in silence the questions
of multiplicative structures, although they are quite essential in applying
the notions of Chapters III, IV, V. Let us note, moreover, that there does not
yet seem to be a satisfactory theory of multiplicative structures in
homological algebra having the necessary degree of generality and simplicity
(Chapter II of [6] being, moreover, a striking illustration of this state of
affairs).\LayerTarget{tohoku:introduction:footnote:1bis}\textsuperscript{1 bis)} For multiplication in cohomology

\TohokuContinuation{121}
